% unit14.tex — R2 unit 14 build driver, The Sourdough Framework: Arabic Edition
%
% Original: The Sourdough Framework by Hendrik Kleinwächter,
% the-sourdough-framework.com, CC BY-SA 4.0. The Arabic edition inherits
% CC BY-SA 4.0.
%
% Unit 14 per the reconstruction plan: ch. 14 troubleshooting, part 1
% (accepted pp. 151–165; sections 14.1–14.3.6 plus the 14.4 section
% head and intro paragraph closing accepted p. 165). Five figures
% (hooch, fungi, ethanol-oxidation chemfig, tearing-dough,
% parbaked-bread), twelve cite keys, one footnote. Content layer
% transcribed from the accepted edition text layer
% (book-ar-fixed.pdf, b78c2a28…, PRIMARY) with the recovery reference
% (recovery-r0/text/14-troubleshooting.txt) as cross-check and
% spacing/wording arbiter — no retranslation; upstream misc.tex gives
% structure only (cite keys, qty values, figure specs, labels). Built
% with the standing arabic-digits.lua (standing md5
% a2e3d1daeb3e24ced21873376170cab8) and run.sh per Ruling 2
% (2026-11-15). Driver scheme per the accepted edition, same as
% units 3–13: NO \frontmatter — chapters are numbered from 1 and page
% numbering runs continuously in Western digits from the cover. Unit
% builds carry natural chapter numbering (ch. 14 builds as ch. 1
% here, figures as 1.x, cites as [1]/[2]…); the R3 full-book assembly
% regenerates true numbering.

\documentclass[paper=a4, twoside=false, fontsize=12pt, parskip=half,
                bibliography=totoc, listof=totoc]{scrbook}
\usepackage{sourdough-ar}
% range-phrase wrapped in \textenglish per the unit 9 dash-font finding
% (2026-11-15), inherited by units 10–14: literal-text ranges
% (\LR{\textenglish{N\textendash M}}) are the accepted convention;
% \numrange with this \textenglish range-phrase is the fallback
% convention. Driver-level change only.
\sisetup{range-phrase={\textenglish{\textendash}}}
\author{Hendrik Kleinw\"achter}
\title{منظومة عمل الخبز الحامض}

\begin{document}
% cover-ar.tex — Arabic cover and title page for The Sourdough Framework
%
% Original: The Sourdough Framework by Hendrik Kleinwächter,
% the-sourdough-framework.com, CC BY-SA 4.0. The Arabic edition inherits
% CC BY-SA 4.0.
%
% Fork of the upstream cover/cover.tex. Page 1 carries the cover raster
% (upstream cover/cover-page.jpg, upstream commit 69636e6) with the Arabic
% overlay of record on the cover band; page 2 is the blank verso; page 3 is
% the typeset title page with the Arabic title, edition date and the
% author/translator credit.
%
% Overlay of record (Phase 26/27, Item 1 applied at R1: 4-word title):
%   Line 1: title, \put(88.45, 7.87), 14/16 Amiri bold, centered
%   Line 2: تأليف: هندريك كلاينفِختر · ترجمة: نائل العيله
%           \put(88.45, 2.92), 8/9 Amiri regular, centered
% The \makeatletter/\if@twoside wrapper is upstream-verbatim (a .tex file;
% the no-\makeatletter rule binds the .sty only).

\thispagestyle{empty}
\setlength{\unitlength}{1mm}
\makeatletter
\if@twoside%
    \noindent\begin{picture}(0,0)(1,-1)
    \put(-16.3,-265){\includegraphics[width=1.33\linewidth]{cover/cover-page.jpg}}
    \end{picture}
\else%
    \noindent\begin{picture}(0,0)(1,-1)
    \put(-23.3,-265){\includegraphics[width=1.33\linewidth]{cover/cover-page.jpg}}
    \end{picture}
\fi%
% Arabic overlay on the cover band (Phase 26/27 of record; 4-word title).
\noindent\begin{picture}(0,0)(1,-1)
    \put(88.45,7.87){\makebox(0,0)[c]{\fontsize{14}{16}\selectfont\bfseries منظومة عمل الخبز الحامض}}%
    \put(88.45,2.92){\makebox(0,0)[c]{\fontsize{8}{9}\selectfont تأليف: هندريك كلاينفِختر · ترجمة: نائل العيله}}%
\end{picture}
\makeatother

\newpage
\thispagestyle{empty}

\rule{1pt}{\textheight} % Vertical line
 % Whitespace between the vertical line and title page text
\hspace{0.05\textwidth}
 % Paragraph box for holding the title page text, adjust the width to move the
% title page left or right on the page
\parbox[b]{0.75\textwidth}{%
{\Huge\bfseries منظومة عمل الخبز الحامض}\\[1\baselineskip] % Title (4-word form, Item 1 of record)
{\large\textenglish{The Sourdough Framework}}\\[2\baselineskip] % Original-title subtitle, accepted p. 2
{\large\textit{الإصدار: 7 نوفمبر 2026}}\\[4\baselineskip] % date frozen to the accepted edition's build date
{\Large هندريك كلاينفِختر \textenglish{(Hendrik Kleinwächter)}}\\[1\baselineskip]
{\large ترجمة: نائل العيله}

% Whitespace between the title block and the copyright text. In the PDF and
% the ebook it holds a note on asking this book: Loafy answers from it. The
% website leaves it out, it has Loafy on every page already. Each place has
% its own via tag, so Loafy can count where readers came from.
\ifdefined\HCode
\vspace{0.5\textheight}
\ifdefined\sdfebook
\textbf{أسئلة أثناء الخَبز؟} اسأل هذا الكتاب. \textenglish{Loafy}، رفيقه المجاني،
يجيب من هذه الصفحات ويسمّي القسم:
\href{https://loafy.ai/loafy?via=tsf-epub-title}{https://loafy.ai}
\fi
\else
\vspace{0.17\textheight}

\noindent
\begin{minipage}[c]{2.4cm}
% Black, not link blue: a light QR code is hard to scan once printed.
% The QR matrix macros must typeset LTR; RTL context corrupts them.
\LRE{\qrcode[nolink,height=2.2cm]{https://loafy.ai/loafy?via=tsf-pdf-qr}}
\end{minipage}%
\hspace{0.4cm}%
\begin{minipage}[c]{\linewidth-2.8cm}
\textbf{أسئلة أثناء الخَبز؟} اسأل هذا الكتاب «لُوفي» \textenglish{(Loafy)}، رفيقه المجاني،
يجيب عنك من صفحاته ويسمّي لك القسم:
\href{https://loafy.ai/loafy?via=tsf-pdf-title}{\nolinkurl{https://loafy.ai}}
\end{minipage}

\vspace{0.17\textheight}
\fi


{\noindent
\begin{flushleft}
    \includegraphics[width=3cm]{cover/CC-BY-SA}\par
الشيفرة المصدرية الكاملة لهذا الكتاب متاحة على
\url{https://github.com/hendricius/the-sourdough-framework/} برخصة
\textenglish{CC-BY-SA}.
راجع \url{https://creativecommons.org/licenses/by-sa/4.0/} لمزيد من التفاصيل.
لا تتردد في الإبلاغ عن الأخطاء أو اقتراح التحسينات. يتوفر إصدار غلاف مقوّى من
الكتاب أيضًا. مزيد من المعلومات هنا:
\url{https://www.breadco.de/hardcover-book}

الإصدار العربي: مبنيّ على الكتاب الأصلي \textenglish{The Sourdough Framework}
لـ هندريك كلاينفِختر، المنشور مرخّصًا برخصة \textenglish{CC BY-SA 4.0}،
ويخضع هذا الإصدار العربي للرخصة نفسها.
\end{flushleft}
}

{\centering\textenglish{Original: The Sourdough Framework by Hendrik
Kleinwächter, the-sourdough-framework.com, CC BY-SA 4.0}\par}
}


{\hypersetup{hidelinks}\tableofcontents}

% troubleshooting-ar.tex — Arabic edition of ch. 14, troubleshooting
%
% Original: The Sourdough Framework by Hendrik Kleinwächter,
% the-sourdough-framework.com, CC BY-SA 4.0. The Arabic edition inherits
% CC BY-SA 4.0.
%
% Phase 28R R2 unit 14 reconstruction (ch. 14 part 1 = accepted
% pp. 151–165, sections 14.1–14.3.6; the 14.4 section head + intro
% paragraph close accepted p. 165 and are included so the last page
% matches 1:1 — unit 15 continues from here with the crumb-structures
% content). METHOD (unit 14 finding): the accepted edition's own text
% layer is PRIMARY for transcription; the R0 recovery reference
% (recovery-r0/text/14-troubleshooting.txt, mined from the Phase 23
% website build of the accepted Arabic sources) is the cross-check and
% the spacing/wording arbiter where the text layer strips spaces; the
% upstream English (book/troubleshooting/misc.tex + crumb-structures.tex)
% gives structure only (figure specs, cite keys, labels, qty values).
% No text retranslated from the English. The accepted edition
% book-ar-fixed.pdf (b78c2a28b574288801d8000a7249af3a) is the fidelity
% authority.
%
% Five figures (upstream [!htb] specs kept verbatim, placement verified
% against the accepted render): hooch 0.5\textwidth (fig:hooch, cite
% liquid+on+starter), fungi-lactic-acid-interactions \textwidth
% (fig:fungi-lactic-acid-interactions, cite + \citeauthor construction
% mirroring the upstream caption), ethanol-oxidation chemfig carried
% over verbatim via \input{figures/fig-ethanol-oxidation.tex} — the
% accepted render of fig 14.3 is Latin-only (verified p. 158),
% tearing-dough 0.8\textwidth (fig:tearing-dough), parbaked-bread
% \textwidth without \centering (upstream verbatim; accepted-edition-
% warranted). LOF short captions recovered from the accepted edition's
% list-of-figures page (accepted p. 195).
%
% One footnote (14.3.2), marker AFTER the sentence period (accepted
% p. 163 token order 'قليلة.1' verified in the render), unlike the
% upstream English (marker before period).
%
% Cites carried with the upstream keys: liquid+on+starter,
% food+safe+ph, mold+lactic+acid+bacteria (cite + \citeauthor),
% acetic+acid+production, rye-defects, oiv+volatile+acidity,
% vriesekoop+spent+grain, purlis+browning, reichenberger+amylase,
% more+active+starter (\cite* verbatim — renders as a plain numeric
% cite in the accepted edition, accepted p. 163 [36]), stiff+starter
% (twice), baking+powder+reduce-acidity.
%
% Direction treatments per the standing conventions:
%  - cross-chapter refs literal (9.2, 9.6, 9.7, 9.9, 6.3, 6.4 —
%    undefined-refs pending set, R3); the 14.4.2 refs are same-chapter
%    but their label lives in unit 15 territory (crumb-structures), so
%    they ride literal \LR{14.4.2} islands for now, pending-set at
%    unit 15 / R3;
%  - the same-unit figure ref renders \LR{\ref{fig:tearing-dough}}
%    (unit build shows 14.4 where the accepted shows 14.4 — unit-
%    numbering artifact, R3 assembly regenerates true numbering);
%  - feeding ratios as LTR islands per the unit 7 precedent:
%    \LR{1:5:5}, \LR{1:10:10}, \LR{1:5:2.5}, \LR{1:1:1};
%  - digit ranges as literal-text LTR islands per the unit 9 dash-font
%    convention: \LR{\textenglish{12\textendash 24}},
%    \LR{\textenglish{2\textendash 3}}, \LR{\textenglish{20\textendash 30}};
%  - bare digit runs ride per-token \LR islands (unit 13 convention):
%    \LR{7} أيام, \LR{24} ساعة, \LR{30} دقيقة/يومًا, \LR{4.2};
%  - \qty/\qtyrange/\numproduct ride \LR{} islands (units 7–13), incl.
%    \LR{\qtyrange{50}{60}{\percent}} and \LR{\qtyrange{1}{2}{\gram}}
%    for the accepted edition's en-dash ranges; temps as
%    \LR{\qty{N}{\degreeCelsius}} with the Fahrenheit parenthesized;
%    \ch{CO2} rides \LR{} per the bread-types precedent;
%  - Latin terms ride \textenglish{} islands: (FAQ), (hooch),
%    (Parmigiano), (ciabatta), (pH), (proteolysis), (baking powder);
%    the \citeauthor output renders in the Latin font by itself (no
%    wrapper needed — accepted p. 156 font census);
%  - the accepted edition's own text-layer tatweel in 'بـ20 g' /
%    'بـ50 g' (pp. 155–157) IS source text here — the recovery
%    reference confirms بـ before the gram islands, matching the
%    bread-types precedent ('وغذِّها بـ\LR{\qty{5}{\gram}}');
%  - verbatim oddities kept as the accepted sources have them (both the
%    text layer and the recovery reference agree): سندرتفع (14.1.1),
%    هي شيوعًا إما (14.1.2), الأصلب (14.1.3), أطرى (14.1.4/14.1.4),
%    يُعرف فيها (14.2.4), قد تخبز (14.2.1), المغذاة (14.2.4), تزايُد
%    (14.2.2), و للأرغفة with the space (14.3.2), وتميل الدقيق /
%    ولهذا تميل (14.3.2), كخَبز (14.3.2), غير كادرة (14.3.1),
%    شأنه شان (14.2.4/14.3.5), ويلعب بلعبٍ (14.1.10), ويسمى الدقيق
%    كخَبز هذا نوعًا (14.3.2); the doubled damma in the accepted text
%    layer's 'أُُجري' (14.2.1 heading) is an extraction artifact — the
%    recovery reference (mined from the accepted sources) has أُجري;
%  - the ``...'' quotes around تهيئة البادئة (14.1.1) are verbatim
%    ASCII backticks/quote-pairs — the accepted render shows them
%    unligated in the Arabic (Amiri) context, accepted p. 152 text
%    layer.

\chapter{استكشاف الأخطاء وإصلاحها}\label{ch:troubleshooting}

\begin{quoting}
يمكنك أن تعتبر هذا الفصل أسئلةً شائعةً \textenglish{(FAQ)} حول أكثر المشكلات التي يواجهها الخبّازون؛ فهو يمنحك أدوات التشخيص التي تحتاجها لتحليل الوضع. ويمكنك بعدها تطبيق الإجراءات المناسبة فتُسقط كل علّة واحدة تلو الأخرى حتى تبلغ الرغيف المثالي.
\end{quoting}

\section{البادئة}

\subsection{بادئتي لا تتضاعف في الحجم}

يدعو بعض الخبّازين إلى انتظار تضاعف البادئة الحامضة في الحجم قبل استخدامها. والفكرة هي استخدام البادئة في ذروة نشاطها لضمان تخمير متوازن في العجين الرئيسي.

غير أن مقياس تضاعف الحجم ينبغي التعامل معه بحذر عند الحكم على بادئتك. فبحسب الدقيق الذي تغذّي به البادئة تتوقّع درجات مختلفة من ارتفاعها. فمثلًا إذا استخدمت دقيق الجاودار فلن تحتفي البادئة سوى بقليلٍ من غازات التخمير داخلها، وبنتيجة ذلك لن ترتفع بادئتك الحامضة كثيرًا، ومع ذلك قد تكون في حالة صحية تامة. وإذا استخدمت دقيق قمحٍ أقل غلوتينًا فلن ترتفع البادئة كثيرًا أيضًا. والسبب أن شبكة الغلوتين لديك أضعف، فتتسرب كميات أكبر من الغاز خارج عجينك.

والمقصود من ذلك أنك مدعوّ إلى أن تنمّي مقياس زيادة الحجم عندك. فبادئتك سندرتفع في الحجم ثم تفقد بنيتها في النهاية وتنكمش. لاحظ اللحظة السابقة للانكماش؛ فهي اللحظة التي ينبغي أن تستخدم فيها بادئتك. وقد تكون هذه زيادةً في الحجم تبلغ \LR{\qty{50}{\percent}} أو \LR{\qty{100}{\percent}} أو \LR{\qty{200}{\percent}}. ومن الأفضل دائمًا أن تستخدم البادئة قبل قليلٍ من أن تكون متأخرًا. فإذا استخدمت البادئة متأخرًا فقلّل الكمية التي تستخدمها. فإذا استدعت الوصفة كمية بادئة تبلغ \LR{\qty{20}{\percent}} فاستخدم في تلك الحالة \LR{\qty{10}{\percent}} فقط. وسيعود نمو بادئتك من جديد داخل عجينك الرئيسي.

وفوق الاعتماد على زيادة الحجم، ابدأ بتدوين ملاحظاتك عن رائحة بادئتك. فمع مرور الوقت ستتمكن من الحكم على حالة تخميرها من الرائحة. وكلما كانت الرائحة أقوى، كان عجينك قد تخمّر أكثر. وهذه إشارة إلى وجوب استخدام بادئة أقل عند صنع العجين الفعلي.

لمزيد من المعلومات حول هذا الموضوع راجع القسم \LR{9.2} ``تهيئة البادئة''.

\subsection{ما أفضل نسبة تغذية للبادئة؟}

أفضل نسبة تغذية للبادئة هي شيوعًا إما \LR{1:5:5} أو \LR{1:10:10}. وفي حالة \LR{1:5:5} تكون قديمة البادئة جزءًا واحدًا، والدقيق خمسة أجزاء، والماء خمسة أجزاء. وإذا كنت تستخدم بادئة صلبة فاستخدم نصف كمية الماء، أي \LR{1:5:2.5}. وبحسب موعد آخر تغذية لبادئتك قد يكون \LR{1:10:10} خيارًا أجدى. فإذا كانت البادئة قديمة ولم تُغذَّ مؤخرًا كانت نسبة \LR{1:10:10} الخيار الأفضل. فبتخفيض نسبة تلقيح البادئة توفّر للكائنات الدقيقة بيئة أنظف، وبهذا تتكاثر وتنمو من جديد في توازن أكثر مرغوبية لبدء تخمير عجينك.

وبصورة عامة، تعامل مع بادئتك الحامضية بوصفها عجينًا. استخدم النسب نفسها التي تستخدمها لعجين خبزك لبادئتك. فينبغي تدريب بادئتك في البيئة نفسها التي ستستخدمها لاحقًا لعجينك. وبهذا تكون بادئتك على أكمل وجه ملائمة لتخمير العجين الذي تُلقَّح به لاحقًا.

والاستثناء الوحيد من قاعدة \LR{1:5:5} و\LR{1:10:10} هو مرحلة إنشاء البادئة الأولية. ففي الأيام الأولى من عملية صنع البادئة لا تتوفر كائنات دقيقة كافية بعد. لذا فإن استخدام نسبة \LR{1:1:1} يمكن أن يسرّع العملية.

\subsection{ما فائدة استخدام بادئة صلبة؟}

البادئة الحامضية العادية تتألف من أجزاء متساوية من الدقيق والماء (ترطيب \LR{\qty{100}{\percent}}). أما البادئة الأصلب فترطيبها يبلغ \LR{\qtyrange{50}{60}{\percent}}.

فالبادئة الصلبة تعزز الجانب الخميري من بادئتك أكثر. وبهذا يتفكك غلوتينك أبطأ ويمكنك التخمير مدةً أطول. وهذا مفيد خصوصًا عند الخَبز بدقيق أقل غلوتينًا.

يمكنك قراءة المزيد حول موضوع البادئات الصلبة في القسم \LR{6.4}.

\subsection{ما فائدة استخدام بادئة سائلة؟}

تدفع البادئة السائلة التخمير اللاهوائي للبكتيريا في بادئتك. وبهذا تميل بادئتك إلى إنتاج حمض اللبنيك أكثر من حمض الخليك. وحمض اللبنيك يُدرَك أطرى وأشبه باللبن الرائب. أما حمض الخليك فقد يكون مذاقه حادًّا أحيانًا. ويمكن أن يكون حمض الخليك مثاليًا عند صنع خبز جاودار داكن، لكنه ليس كذلك كثيرًا عند صنع رغيف منفّاف على نمط تشباتا \textenglish{(ciabatta)}.

وعند تحويل بادئتك إلى بادئة سائلة فأنت تغيّر البيئة الميكروبية لبادئتك تغييرًا دائمًا. ولا يمكنك التراجع بعد أن تقضي على بكتيريا إنتاج حمض الخليك. لذا يُنصح بالاحتفاظ بنسخة احتياطية من بادئتك الأصلية.

ومن سلبيات البادئة السائلة ارتفاع النشاط البكتيري كليًّا مقارنةً بالنشاط الخميري. ومعنى ذلك أن الخبز المخبوز سيكون أكثر حموضة (لكن أطرى). وسيتفكك العجين أسرع أثناء التخمير. ولهذا السبب ستحتاج إلى استخدام دقيق قوي عالي الغلوتين عند استخدام هذا النوع من البادئات.

يمكنك قراءة المزيد عن البادئة السائلة في القسم \LR{6.3}.

\subsection{بادئتي الجديدة لا ترتفع إطلاقًا}

تأكد من استخدامك ماءً غير مكلور. ففي مناطق كثيرة من العالم يُضاف إلى ماء الصنبور كلور لقتل الكائنات الدقيقة. وإذا كان الأمر كذلك في منطقتك فسيساعدك ماء ينبوع معبّأ. ويمكنك أيضًا استخدام مرشّح ماء بفحم نشط سيزيل الكلور. أو بدلًا من ذلك، إذا صببت ماء الصنبور في إبريق أو وعاء آخر وتركيته مغطىً بتغطية فضفاضة، فينبغي أن يتبخر الكلور خلال \LR{\textenglish{12\textendash 24}} ساعة، وتكسب ميزة إضافية وهي امتلاكك تلقائيًا ماءً بحرارة الغرفة.

احرص على استخدام دقيق كامل الحبة (قمحًا كاملًا أو جاودارًا كاملًا، إلخ). فلحبوب الدقيق هذه تلوّث طبيعي أكبر بالخمائر البرية والبكتيريا. وصنع بادئة من الدقيق الأبيض وحده أحيانًا لا ينجح. حاول استخدام دقيق عضوي غير مبيّض لصنع البادئة. فالدقيق الصناعي يُعالَج أحيانًا بمبيدات الفطريات.

\subsection{صنعتُ بادئة، ارتفعت في اليوم الثالث ثم لم تعد ترتفع}

هذا أمر طبيعي. فبما أن بادئتك تنضج تُستثار كائنات دقيقة مختلفة. وخصوصًا خلال الأيام الأولى من العملية يمكن أن تُستثار كائنات دقيقة ضارة مثل العَفَن. وهذه تجعل بادئتك ترتفع كثيرًا. ومع كل تغذية لاحقة للبادئة تنتقي الكائنات الدقيقة الأكفأ في تخمير الدقيق. ولهذا السبب يُنصح بالتخلص من البادئة الزائدة غير المستخدمة من الأيام الأولى للعملية. ولاحقًا لا ينبغي إلقاء كميات البادئة غير اللازمة أبدًا. يمكنك صنع خبز رائع من البادئة الزائدة.

فواصل العملية ولا تستسلم. فالارتفاع الأول الكبير مؤشر على أنك تفعل كل شيء على نحو صحيح. وبخبرتي يستغرق نمو البادئة نحو \LR{7} أيام. وكلما غذّيت بادئتك أكثر وأكثر أصبحت أكفأ في تخمير الدقيق. وقد لا يخرج خبزك الأول تمامًا كما خططت، لكنك ستصل إلى هناك في النهاية. وكل تغذية تجعل بادئتك أقوى فأقوى.

\subsection{سائل فوق بادئتي}

أحيانًا يتجمّع سائل، وفي كثير من الحالات سائل أسود، فوق بادئتك الحامضية. وقد تكون للسائل رائحة نفّاذة. وكثير من الناس يخلطون بين هذا والعَفَن. ورأيتُ خبّازين ينصحون بالتخلص من البادئة بسبب هذا السائل. ويُعرف هذا السائل عادةً بالهُوش \textenglish{(hooch)}. فبعد فترة من انعدام النشاط يفصل الدقيق الأثقل نفسه عن الماء. فيجلس الدقيق في قاع جرّتك ويبقى السائل فوقه. ويسودّ السائل لأن بعض جسيمات الدقيق أخف وزنًا من الماء فتطفو فوقه. وعلاوةً على ذلك تطفو كائنات دقيقة ميتة في هذا السائل. وليس هذا السائل شيئًا سيئًا؛ فهو يحمي بادئتك الحامضية بحمايةٍ فعّالة من دخول العَفَن الهوائي من الأعلى.

\begin{figure}[!htb]
  \centering
  \includegraphics[width=0.5\textwidth]{sourdough-starter-hooch}
  \caption[الهُوش] {هُوش \textenglish{(hooch)} يتجمّع فوق بادئة حامضية~\cite{liquid+on+starter}.}%
  \label{fig:hooch}
\end{figure}

فما عليك إلا أن تحرّك بادئتك الحامضية لتخلط الهُوش من جديد داخل بادئتك. وسيختفي الهُوش. ثم استخدم قليلًا من بادئتك الحامضية لإعداد البادئة الخاصة بخبزك القادم. وبمجرد ظهور الهُوش فبادئتك قد تخمّرت على الأرجح فترةً طويلة. وقد تكون حامضةً جدًّا. وهذه حالة البادئة ممتازة لصنع بقسماط من البادئة الزائدة أو خبز من البادئة الزائدة. لا ترمِ شيئًا. فالهُوش إشارة إلى أن عجينًا طويل التخمير أمامك. قارنه بجبن بارميزان \textenglish{(Parmigiano)} ناضج لسنتين. فالعجين الذي أمامك مليء بالنكهة الشهية.

\subsection{معالجة بادئة متعفّنة}

أولًا وقبل كل شيء، صنع بادئة حامضية متعفّنة أمر صعب للغاية. فهو مؤشر على أن شيئًا ما ربما اختلّ تمامًا في بادئتك. فالتعايش الطبيعي بين الخميرة والبكتيريا لا يسمح عادةً لمسبّبات أمراض خارجية مثل العَفَن بدخول بادئتك الحامضية. والحموضة المنخفضة التي تخلقها البكتيريا بيئة معادية جدًّا لا يحبها أيٌّ من مسبّبات الأمراض الأخرى. وبصورة عامة كل ما دون درجة حموضة تبلغ 4.2 يمكن اعتباره آمنًا غذائيًا~\cite{food+safe+ph}. وهذه هي فكرة الأطعمة المخلّلة. وخبزك الحامض في جوهره خبزٌ مخلّل.

ورأيتُ هذا يحدث خصوصًا عندما تكون البادئة الحامضية صغيرة السنّ نسبيًا. فكل دقيق يحتوي طبيعيًا على أبواغ عَفَن. وعند بدء صنع بادئة حامضية تبدأ جميع الكائنات الدقيقة بالتنافس في استقلاب الدقيق. وقد يفوز العَفَن بالسباق أحيانًا ويغلب الخمائر البرية والبكتيريا الطبيعية. وفي تلك الحالة جرّب ببساطة ترويض بادئتك الحامضية من جديد. وإذا ظهر العَفَن مرة أخرى فقد يكون ذلك دفعةَ دقيقٍ كثيرِ العَفَن. جرّب دقيقًا مختلفًا لبدء بادئتك الحامضية به.

فالبادئات الحامضية الناضجة لا ينبغي أن تتعفّن ما لم تتغيّر ظروف البادئة. ورأيتُ العَفَن يظهر عندما تُخزَّن البادئة في الثلاجة ويجفّ سطحها. ويتكوّن أحيانًا أيضًا على حوافّ وعاء بادئتك، عادةً في المناطق التي لا تبلغها كائنات بادئة نشطة. فجرّب استخراج منطقة من بادئتك لا عَفَن فيها. وغذّها من جديد بالدقيق والماء. وبعد بعض التغذيات ينبغي أن تعود بادئتك إلى طبيعتها. خذ قليلًا جدًّا من البادئة: \LR{\qtyrange{1}{2}{\gram}} تكفي. فهي تحتوي أصلًا على ملايين الكائنات الدقيقة.

فالعَفَن يفضّل الظروف الهوائية. ومعنى ذلك أن الهواء لازم لكي ينمو فطر العَفَن. وتقنية أخرى نحت معي كانت تحويل بادئتي الحامضية إلى بادئة سائلة. وقد نجح هذا في تحويل بادئتي من إنتاج حمض الخليك إلى إنتاج حمض اللبنيك. فحمض الخليك، شأنه شان العَفَن، يحتاج إلى أكسجين ليُنتَج. وبعد غمر الدقيق بالماء تغلّبت بكتيريا حمض اللبنيك بمرور الوقت على بكتيريا حمض الخليك. وهذا مفهوم شبيه بالأطعمة المخلّلة. وبفعل ذلك أنت تقتل عمليًا كل فطريات العَفَن الحية. وقد لا يبقى عندك سوى بعض الأبواغ. ومع كل تغذية ستقلّ الأبواغ قليلًا فقليلًا. وعلاوةً على ذلك يبدو أن بكتيريا حمض اللبنيك تنتج أيضًا استقلابية تثبّط نمو العَفَن~\cite{mold+lactic+acid+bacteria}.

\begin{figure}[!htb]
  \centering
  \includegraphics[width=\textwidth]{fungi-lactic-acid-interactions}
  \caption[تفاعل بكتيريا حمض اللبنيك وفطريات العَفَن] {تفاعل بكتيريا حمض اللبنيك وفطريات العَفَن. في~\cite{mold+lactic+acid+bacteria} يعرض \citeauthor{mold+lactic+acid+bacteria} وآخرون كيف تنتج البكتيريا أيضًا استقلابية تثبّط نمو الفطر.}%
  \label{fig:fungi-lactic-acid-interactions}
\end{figure}

ولتخليل بادئتك خذ ببساطة قليلًا من بادئتك القائمة (\LR{\qty{5}{\gram}} مثلًا). ثم غذّ الخليط بـ\LR{\qty{20}{\gram}} من الدقيق و\LR{\qty{100}{\gram}} من الماء. تكون بذلك قد صنعت بادئة ترطيبها نحو \LR{\qty{500}{\percent}}. رجّ الخليط رجًّا قويًا. وبعد بضع ساعات ينبغي أن تبدأ في رؤية معظم الدقيق قرب قاع وعائك. وبعد حين يستنفد معظم الأكسجين من خليط القاع وتبدأ بكتيريا حمض اللبنيك اللاهوائية في الازدهار. دوّن ملاحظةً عن رائحة بادئتك الحامضية. فإذا كانت خِلّيةً سابقًا فستتغيّر لتصبح لبنيةً كثيرًا. استخرج قليلًا من خليطك في اليوم التالي برجّ كل شيء أولًا. خذ \LR{\qty{5}{\gram}} من الخليط السابق، وغذّ من جديد بـ\LR{\qty{20}{\gram}} آخر من الدقيق و\LR{\qty{100}{\gram}} آخر من الماء. وبعد \LR{\textenglish{2\textendash 3}} تغذيات إضافية ينبغي أن تكون بادئتك قد تأقلمت. وعند العودة إلى ترطيب يبلغ \LR{\qty{100}{\percent}} ينبغي أن يكون العَفَن قد استأصل. لاحظ رجاءً أن إجراء المزيد من الاختبارات حول هذا الموضوع واجب. وسيكون من الجميل تحليل الكائنات الدقيقة بعناية فائقة قبل التخليل وبعده.

\subsection{بادئتي حامضة أكثر من اللازم}

إذا كانت بادئتك الحامضية حامضة أكثر من اللازم فستسبب مشكلات أثناء التخمير. فسيكون في تخميرك نشاط بكتيري أكثر من النشاط الخميري. ومعنى ذلك أنك ستصنع على الأرجح رغيفًا أكثر حموضة وأقل منفّافية مما يمكن أن يكون. والهدف هو بلوغ التوازن الصحيح: قوام منفّاف من الخميرة، وحموضة لاذعة عظيمة ليست بالقوية المفرطة من البكتيريا. وهذا يتوقف طبعًا على ما تبحث عنه من مذاق في خبزك. فعند صنع خبز الجاودار أفضّل مثلًا أن أكون في الجهة الحمضية اللاذعة. وعندما يختل التوازن الموصوف فإن أول ما تفحصه هو بادئتك الحامضية.

لاحظ رائحة بادئتك. هل تنبعث منها رائحة حامضة جدًّا؟ ذُق قليلًا من بادئتك أيضًا. كم هي حامضة المذاق؟ مع مرور الوقت تصبح كل بادئة أكثر حموضةً فأكثر كلما انتظرت. لكن أحيانًا تصبح بادئتك حامضة بسرعة كبيرة. وفي هذه الحالة طبّق تغذيات يومية على بادئتك. وقلّل كمية البادئة القديمة التي تستخدمها للتغذية. ونسبة \LR{1:5:5} أو \LR{1:10:10} قد تفعل العجائب. وفي هذه الحالة ستأخذ جزءًا واحدًا من البادئة (\LR{\qty{10}{\gram}}) وتغذّيه بـ\LR{\qty{50}{\gram}} من الدقيق و\LR{\qty{50}{\gram}} من الماء. وبهذا تبدأ الكائنات الدقيقة التخمير في بيئة بكر جديدة. وهذا شبيه بنسبة بادئة \LR{\qty{10}{\percent}} أو بادئة \LR{\qty{20}{\percent}} التي تستخدمها لصنع عجين. هذه الأيام لم أستخدم نسبة \LR{1:1:1} تقريبًا أبدًا. فهي منطقية فقط عندما تكون في البداية تصنع بادئتك. تريد بيئة حامضة لكي تغلب كائناتك الدقيقة مسبّبات الأمراض المحتملة. فالبيئة الحمضية سامة لمعظم مسبّبات الأمراض التي لا تريدها في بادئتك.

وأسلوب آخر قد يساعد هو تحويل بادئتك الحامضية إلى بادئة صلبة كما ورد في القسم \LR{6.4}.

\subsection{لماذا تفوح من بادئتي رائحة الخل أو الأسيتون؟}

على الأرجح أن بادئتك الحامضية أنتجت كمية كبيرة من حمض الخليك. وحمض الخليك أساسي عند صنع الخل. وحين لا يبقى طعام إضافي تستهلك بعض بكتيريا بادئتك الإيثانول وتحوّله إلى حمض الخليك. وحمض الخليك له رائحة نفّاذة جدًّا. وعند تذوّق حمض الخليك كثيرًا ما يُدرَك مذاق خبزك على أنه قوي جدًّا.

\begin{figure}[!htb]
  \centering
  % chemfig laid out LTR per the accepted render of fig 14.3 (accepted
  % p. 158: ethanol left, arrow right, acetic acid right): bare \input
  % under the RTL paragraph direction mirrors the whole equation and
  % spreads it to 555.8pt (accepted: 329.5pt). \LR keeps the chemistry
  % LTR; the caption stays RTL.
  \LR{\input{figures/fig-ethanol-oxidation.tex}}
  \caption[صنع حمض الخليك] {الأكسجين لازم لصنع حمض الخليك~\cite{acetic+acid+production}.}%
  \label{fig:ethanol-oxidation}
\end{figure}

وليس هذا أمرًا سيئًا. لكن إذا أردت تغيير مذاق خبزك النهائي ففكر في تحويل بادئتك الحامضية إلى بادئة سائلة. سيساعد هذا في إعطاء الأولوية لبكتيريا إنتاج حمض اللبنيك. سيتغيّر مذاقك من خِلّي إلى لبنّي. ولا يمكنك العودة خلفًا مع ذلك. فبعد التحويل لن تعود بادئتك أبدًا إلى إنتاج حمض الخليك لأنك غيّرت تيارات الميزان نحو التخمير اللبني في المقام الأول. أحبّ أن أمتلك بادئة جاودار منفصلة. وفي تجاربي تميل بادئات الجاودار إلى توفر الكثير من بكتيريا حمض الخليك. وهذه البادئة ممتازة عندما تريد صنع خبز قوي الطعم كثير الديمومة. فخبز الجاودار الخالص يذوق ممتازًا إذا صُنع بمثل هذه البادئة. فالمذاق عند القضم لا يوصف. ويلعب بلعبٍ جميل مع الحساء أيضًا. خذ فقط قليلًا من هذا الخبز واغمسه في حسائك.

\subsection{لماذا لا تطفو بادئتي بعد إجراء اختبار الطفو؟}

قد لا يحدد اختبار الطفو جاهزية بادئتك للتلقيح في عجينك تحديدًا موثوقًا. فهو فعّال في عجائن القمح حيث تُحتجَز كمية وفيرة من الغاز داخل شبكة الغلوتين، لكنه أقل موثوقية في العجائن غير القمحية. ففي العجائن غير القمحية يميل الغاز المتولد أثناء التخمير إلى الهروب، فتغرق البادئة على الأرجح.

وللحكم على جاهزية بادئتك بدقة أكبر راقب الفقاعات عند حافّة الوعاء وخذ رائحتها في الحسبان. فالبادئة الناضجة ينبغي أن تنبعث منها رائحة حامضة معدلة دون أن تكون نفّاذة أكثر من اللازم.

\section{العجين}

\subsection{هل أُجري أوتولايز لعجيني؟}

في \LR{\qty{95}{\percent}} من جميع الحالات لا معنى لأوتولايز. وبدلًا من ذلك أنصحك بإجراء الفيرمنتوليز. يمكنك قراءة المزيد عن عملية أوتولايز في القسم \LR{9.6} والمزيد حول موضوع الفيرمنتوليز في القسم \LR{9.7}.

يجمع الفيرمنتوليز كل فوائد أوتولايز مع استبعاد سلبيات مثل اضطرارك إلى عجن العجين أكثر من مرة.

لا معنى لأوتولايز إلا إذا كنت قد تخبز عجينًا سريع التخمير قائمًا على الخميرة وبمعدل تلقيح خميري مرتفع. لكن حتى في تلك الحالة يمكنك ببساطة تخفيض كمية الخميرة لتُجري الفيرمنتوليز بدل أوتولايز.

% Heading parens wrapped in \textenglish per the binding unit 11
% sans-context Latin-routing convention (bare ASCII in a sans heading
% context routes to the Latin sans face; the accepted edition sets
% these parens in OpenSans-Bold): keeps the text layer clean (bare
% parens in NotoNaskhArabic-Bold map to U+0000).
\subsection{عيّنة عجيني\textenglish{(}جرّة أليكوت\textenglish{)} لا ترتفع. ما الخطأ؟}

إذا رأيت أن عجينك يرتفع في الحجم لكن جرّتك لا ترتفع، فالأرجح أن كليهما يتخمّر بسرعتين مختلفتين. وقد يحدث هذا كثيرًا عندما تتغير الحرارة في مطبخك. فالجرّة أشد تأثرًا بتغيّرات الحرارة من العجين الرئيسي. ولأن حجم العيّنة أصغر فهي تسخن أو تبرد أسرع.

ولهذا السبب يوجب عليك استخدام ماء بحرارة الغرفة عند صنع عجينك. فبالتماثل في الحرارة بين العيّنة وعجينك تضمن أن كليهما يتخمّر بالسرعة نفسها.

وإذا تغيّرت حرارة غرفتك تغيّرًا كبيرًا خلال النهار فخيارك الأفضل هو استخدام وعاء شفّاف. ضع على الوعاء علامة لقياس زيادة حجم عجينك قياسًا سليمًا.

وخيار آخر قد يكون استخدام مقياس الحموضة \textenglish{(pH)} الأكثر تكلفةً لقياس تزايُد حموضة عجينك. يمكنك قراءة المزيد عن الطرق المختلفة لإدارة التخمير الأول في القسم \LR{9.9}.

% \textenglish paren wrappers per the same convention (see above).
\subsection{ما مستوى الماء الجيد\textenglish{(}الترطيب\textenglish{)} لصنع عجين؟}

خصوصًا عندما تبدأ في صنع الخبز استخدم كميات ماء أقل. فسيبسّط هذا العملية كلها تبسيطًا كبيرًا. أنصح باستخدام مستوى ترطيب يبلغ نحو \LR{\qty{60}{\percent}}. فلكل \LR{\qty{100}{\gram}} من الدقيق استخدم نحو \LR{\qty{60}{\gram}} من الماء. وهذا الرقم التقريبي سيناسب معظم الدقيق. وبهذا الترطيب يمكنك صنع خبز ولفائف وبيتزا وحتى باجيت من العجين نفسه.

ومع الترطيب المنخفض يسهل التعامل مع العجين وتزيد نسبة التخمير الخميري، وينتج عن ذلك انخفاض خطر التخمير الزائد.

\subsection{عجيني يتمزق تمامًا بعد تخمير طويل}

أحيانًا عند ملامسة عجينك بعد تخمير طويل يتمزق تمامًا. وقد يكون لذلك سببان. فقد يكون أن البكتيريا استهلكت غلوتين دقيقك بالكامل. ومن جهة أخرى تتفكك شبكة غلوتينك تلقائيًا مع مرور الوقت. وهذه عملية يُعرف فيها بالإنزيم البروتييز الذي يحوّل شبكة الغلوتين إلى أحماض أمينية أصغر يمكن للبذرة استخدامها لبناءً لنموها. تُعرف هذه العملية باسم تفكك البروتين \textenglish{(proteolysis)}، وتبدأ لحظة خلطك الدقيق بالماء. وكلما جلس عجينك مدةً أطول تكسّر غلوتين أكثر. وبما أن الغلوتين هو ما يمسك عجين القمح متّصلًا، فسينفطر عجينك في النهاية.

\begin{figure}[!htb]
  \centering
  \includegraphics[width=0.8\textwidth]{tearing-dough}
  \caption[تمزّق العجين] {عجيني يتمزّق بعد \LR{24} ساعة من انعدام النشاط.}%
  \label{fig:tearing-dough}
\end{figure}

في الصورة \LR{\ref{fig:tearing-dough}} جرّبتُ استخدام بادئة لم تُغذَّ لمدة \LR{30} يومًا بحرارة الغرفة. حاولتُ صنع عجين مباشرة من البادئة غير المغذاة. وعادةً بعد فترة طويلة بلا تغذيات تبدأ كائناتك الدقيقة في التبويغ وتدخل طور السبات. وبهذا تنجو فترةً طويلة بلا تغذيات إضافية. وإضافة طعام إضافي ستعيد تنشيطها. وفي هذه الحالة لم يتخمّر العجين بسرعة كافية قبل أن يفكّك البروتييز الغلوتين. وبتنشيط كائناتك الدقيقة ستبدأ بالتكاثر وزيادة العدد ما دام طعامًا متاحًا. لكن هذه العملية في حالتي لم تكن سريعة بما يكفي. وبعد نحو \LR{24} ساعة بدأ العجين كله يتمزق تمامًا. وسرّعتُ العملية كلها باستخدامي دقيق القمح الكامل. فالقمح الكامل يحتوي إنزيمات أكثر من الدقيق الأبيض.

ولمعالجة ذلك تأكد من أن بادئتك الحامضية مرِحة ونشطة. ما عليك إلا أن تزيد التغذيات قليلًا قبل صنع عجينك. وبهذا يصبح عجينك جاهزًا للتشكيل قبل أن يتفكك تمامًا.

\section{الخبز}

\subsection{خبزي يبقى مسطحًا}

الخبز المسطح مرتبط في معظم الحالات بتفكك شبكة الغلوتين لديك تفككًا تامًّا. وليس هذا سيئًا؛ معناه أنك تأكل طعامًا مخمّرًا تمامًا. لكن من منظور المذاق والقوام قد يكون خبزك حامضًا أكثر من اللازم أو لم يعد منفّافًا. ولاحظ رجاءً أيضًا أنك لا تستطيع صنع خبز بانتفاخٍ عظيم في الفرن إلا عند صنع عجائن قمحية. فعندما بدأت بهذه الهواية كنت أتساءل دائمًا لماذا تخرج أخباز الجاودار لدي مسطحة إلى هذا الحد. فالجاودار فيه غلوتين فعلًا، لكن فيه جسيمات صغيرة تُدعى البنتوزانات (أرابينوكسيلان وبيتا-جلوكان)~\cite{rye-defects}. وهي تمنع العجين من تطوير شبكة غلوتين بالقدر نفسه الذي يطوره مع القمح. جهودك ستكون عبثًا، وسيبقى عجينك مسطحًا. فالسبلت والقمح وحدهما قادران على احتجاز \LR{\ch{CO2}} المتولد من التخمير.

وفي معظم الحالات فالأرجح أن شيئًا ما مختل في بادئتك الحامضية. ويحدث هذا كثيرًا جدًّا عندما تكون البادئة صغيرة السنّ نسبيًا وغير كادرة على تخمير الدقيق بعد. ومع مرور الوقت ستصبح بادئتك الحامضية أفضل فأفضل. احتفظ ببادئتك الحامضية بحرارة الغرفة ثم طبّق تغذيات يومية بنسبة \LR{1:5:5}. أي جزء واحد قديم البادئة، وخمسة أجزاء دقيق، وخمسة أجزاء ماء. يتيح لك هذا بلوغ توازن أفضل بين الخميرة والبكتيريا في بادئتك الحامضية. وقد يكون أفضل من ذلك استخدام بادئة صلبة. فالبادئة الصلبة تعزز الجانب الخميري من بادئتك. ويسمح لك هذا بجعل التخمير البكتيري أقل، فينتج عن ذلك شبكة غلوتين أقوى قرب نهاية التخمير~\cite{stiff+starter}.

وراجع رجاءً أيضًا القسم الفرعي \LR{14.4.2} حيث شرحتُ المزيد عن العجائن المتخمّرة زيادةً. ويمكنك أيضًا الرجوع إلى القسم \LR{6.4} لمزيد من التفاصيل عن صنع بادئة صلبة.

وعلاوةً على ذلك، فالدقيق الأقوى الحاوي على غلوتين أكثر سيساعدك على دفع التخمير أبعد. والسبب أن دقيقك يحتوي غلوتينًا أكثر وسيستغرق تفككه بواسطة بكتيرياك مدةً أطول. وفي النهاية، إذا خُمّر مدةً طويلة زائدة، فسيتفكك عجينك أيضًا وسيصبح لزجًا ومسطحًا.

ولتشخيص ما إذا كان التخمير البكتيري الزائد هو المشكلة، ذُق عجينك ببساطة. هل مذاقه حامض جدًّا؟ إذا كان كذلك فهذا مؤشر جيد. وعند العمل بالعجين، هل صار فجأة لزجًا جدًّا بعد بضع ساعات؟ فذلك مؤشر جيد آخر. واستخدم أنفك أيضًا لتدوين ملاحظاتك عن رائحة العجين. فلا ينبغي أن تكون نفّاذة أكثر من اللازم.

\subsection{أريد حموضة أكثر في خبزي}

لبلوغ حموضة أكبر في خبزك الحامض يوجب عليك تخمير عجينك فترةً أطول. ومع مرور الوقت ستستقلب البكتيريا معظم الإيثانول الذي صنعته الخميرة في عجينك. وتنتج البكتيريا في معظمها حمضًا لبننيًا وحمض الخليك. وحمض اللبنيك أكثر حموضةً كيميائيًا من حمض الخليك لكنه أحيانًا لا يُدرَك حامضًا. وفي معظم الحالات فالتخمير الأطول هو ما تريده. وستحتاج إما إلى استخدام قالب الخبز لصنع عجينك أو إلى استخدام دقيق يقاوم فترة تخمير طويلة. ويسمى الدقيق كخَبز هذا نوعًا من الدقيق القوي. وتميل الدقيق الأقوى إلى أن تكون من أصناف قمح نبتت في ظروف أكثر مشمسًا. ولهذا تميل الدقيق الأقوى إلى أن تكون أكثر تكلفة. و للأرغفة القائمة بذاتها أنصح باستخدام دقيق يحتوي \LR{\qty{12}{\percent}} من البروتين على الأقل. وبصورة عامة، كلما زاد البروتين أمكنك تخمير عجينك مدةً أطول.

وخيار آخر لبلوغ نكهة أكثر حموضة قد يكون استخدام بادئة تنتج حمض الخليك أكثر. وبخبرتي الخاصة أنتجت معظم بادئاتي من الجاودار الخالص نكهات خِلّية أقوى. وكيميائيًا ليس حمض الخليك حامضًا بذلك الحد، لكنه عند التذوق يبدو أكثر حموضة. احرص على استخدام بادئة ترطيبها نحو \LR{\qty{100}{\percent}}. فإنتاج حمض الخليك يقتضي أكسجين. وتميل البادئة شديدة السائلة إلى تفضيل إنتاج حمض اللبنيك لأن الدقيق مغمور بالماء. وبغمر العجين لا يمكن سوى لقليلٍ جدًّا من الأكسجين أن يمر عبر الماء إلى الدقيق المتخمّر. ولهذا لا يمكن إنتاج سوى قليلٍ جدًّا من حمض الخليك. ومع مرور الوقت ستؤول بكتيريا إنتاج حمض الخليك من بادئتك فناءً.

\begin{figure}[!htb]
  \includegraphics[width=\textwidth]{parbaked-bread}
  \caption[خبز شبه مخبوز] {خبز شبه مخبوز \textenglish{(parbaked)}.}%
  \label{fig:parbaked-bread}
\end{figure}

وخيار آخر أسهل قد يكون خبز البادئة الحامضية مرتين. وقد لاحظتُ هذا عند شحن الخبز لمخبزتي المصغّرة. كانت الفكرة أن أخبز خبزي نحو \LR{30} دقيقة حتى يُستهلّ، ثم أتركه يبرد ثم أشحنه إلى الزبائن. وبمجرد استلامك له تخبزه مجددًا لنحو \LR{\textenglish{20\textendash 30}} دقيقة أخرى لبلوغ القشرة المطلوبة ثم يمكنك أكله. وأفاد بعض الزبائن بأن الخبز ذو مذاق حامض جدًّا. وأفضل تفسير لدي هو أن خبز الخبز مرتين يُفقدك حموضة أقل أثناء عملية الخَبز. فحمض الخليك يغادر خبزك مع بخار الماء، ومعظمه من القشرة، بينما حمض اللبنيك بالكاد يتبخر~\cite{oiv+volatile+acidity,vriesekoop+spent+grain}. فبعد الخَبز \LR{30} دقيقة عند نحو \LR{\qty{230}{\degreeCelsius}} (\LR{\qty{446}{\degF}}) ما تزال القشرة شاحبة ولم يغادر سوى بعض الماء وحمض الخليك. وإذا واصلت الخَبز فسيمضي حمض الخليك أكثر فأكثر مع البخار. والآن إذا أوقفت الخَبز بعد \LR{30} دقيقة فستكون قد بلغت على الأرجح حرارةً قلبية تبلغ نحو \LR{\qty{95}{\degreeCelsius}} (\LR{\qty{203}{\degF}}). سيتوجب على عجينك أن يبرد مجددًا إلى حرارة الغرفة. وستبقى القشرة شاحبة بقدرٍ ما. ثم بعد بضع ساعات تبدأ في خبز عجينك من جديد. وستصبح قشرتك داكنة جميلة ذات عِطر شهي. والعِطر قادم من تفاعل مايار. فقلب عجينك لا يسخن قط أبعد من نحو \LR{\qty{100}{\degreeCelsius}} (\LR{\qty{212}{\degF}}) ما دام يحتوي ماءً، فاللُّبّ يحتفظ بمعظم حموضه~\cite{purlis+browning,reichenberger+amylase}. وبالمحصّلة سيكون خبزك أكثر حموضة. ويساعد هذه الحموضة المعززة أيضًا في منع مسبّبات الأمراض من دخول خبزك. وسيكون الخبز جيدًا فترةً أطول. ولهذا السبب يعمل مفهوم المخبز التوصيلي جيدًا مع الخبز الحامض اللاذع. وفي تجاربي الخاصة بقي الخبز جيدًا حتى أسبوع في كيس بلاستيكي. وهذا أطول بكثير من عجين قائم على الخميرة قد يُعدّيه العَفَن بعد أيام قليلة.\footnote{غير أن بعض أوائل زبائني التجريبيين أفادوا بأن الخبز كان حامضًا إلى حدٍّ بعيد، غير ممتع به إطلاقًا. وإذا حدث لك هذا ففكر في تحميص الخبز. فالتحميص يطرد حموضةً إضافية بالغليان.}

\subsection{خبزي حامض أكثر من اللازم}

يحب بعض الناس الخبز الأقل حموضةً أيضًا. وهذا تفضيل شخصي. ولبلوغ خبز أقل حموضة يوجب عليك التخمير فترةً أقصر. تنتج الخميرة \LR{\ch{CO2}} والإيثانول. وتستهلك الخميرة والبكتيريا معًا السكريات التي أطلقها إنزيم الأميلاز في عجينك. وعندما يستنفد السكر تبدأ البكتيريا في استهلاك الإيثانول المتبقي من الخميرة. ومع مرور الوقت تُخلَق حموضة أكثر فأكثر، فيصنع ذلك رغيفًا أكثر حموضة.

وزاوية أخرى لهذا قد تكون تغيير نسبة الخميرة/البكتيريا في بادئتك الحامضية. يمكنك أن تبدأ التخمير بخميرة أكثر وبكتيريا أقل. وبهذا، من أجل الزيادة نفسها المعطاة في حجم عجينك، ستكون لديك حموضة أقل. وحيلة جيدة جدًّا هي أن تضمن أنك تغذّي بادئتك مرة في اليوم بحرارة الغرفة. وبهذا تدير تيارات بادئتك نحو تخمير خميري أفضل~\cite*{more+active+starter}.

ولإدارة التيارات أبعد، فقد كان صنع بادئة صلبة عندّي تغييرًا حقيقيًا في قواعد اللعبة. فالبادئة الصلبة ترطيبها نحو \LR{\qty{50}{\percent}}. وبفعل ذلك ستفضّل بادئتك الحامضية النشاط الخميري كثيرًا أكثر. وستكون عجائنك أكثر نفّافية وأقل حموضة من أجل زيادة الحجم نفسها. اختبرتُ ذلك بوضع بالونات فوق جرار زجاجية مختلفة. استخدمتُ الكمية نفسها من الدقيق لكل عيّنة. اختبرتُ بادئة عادية وبادئة سائلة وبادئة صلبة. وخلقت البادئة الصلبة \LR{\ch{CO2}} بكثير أكثر من البادئات الأخرى. ونتيجة لذلك انتفخ بالون البادئة الصلبة الانتفاخ الأكثر~\cite{stiff+starter}. يمكنك قراءة المزيد حول موضوع البادئات الصلبة في القسم \LR{6.4}.

وأنصح نهجًا غير تقليدي آخر قد يكون إضافة مسحوق الخبز \textenglish{(baking powder)} إلى عجينك. فمسحوق الخبز يعادل حمض اللبنيك ويجعل العجين ألطف كثيرًا~\cite{baking+powder+reduce-acidity}.

\subsection{خبزي ينبطح عند إخراجه من سلة التخمير}

بعد إخراج عجينك من سلة التخمير سينبطح عجينك قليلًا دائمًا. والسبب أن شبكة الغلوتين لديك تسترخي مع الوقت ولم تعد تستطيع الإمساك بالشكل. ومع ذلك، خلال سياق الخَبز، سيرتفع عجينك في الحجم وينفخ من جديد.

ولكن إذا انبطح عجينك تمامًا فتلك إشارة إلى أنك خمّرت عجينك مدةً أطول من اللازم. راجع القسم الفرعي \LR{14.4.2} حيث أشرح عن العجائن المتخمّرة زيادةً. فقد استهلكت بكتيرياك معظم شبكة الغلوتين لديك. ولهذا السبب ينهار عجينك تمامًا ويبقى مسطحًا أثناء الخَبز. وسينفذ \LR{\ch{CO2}} والماء المتصاعد خارج العجين. وعرَضٌ شبيه أن عجينك يلتصق بسلة التخمير. وعندما بدأت الخَبز جاهدتُ هذا بدقيق الأرز. وقد نحلت معي لكن قد يكون اكتشافًا زائفًا. راجع القسم الفرعي \LR{14.4.2} لمزيد من التفاصيل عن سبب عدم كون دقيق الأرز فكرةً جيدة لإدارة العجائن اللاصقة.

هذه الأيام أفرك عجيني بلطفٍ بقليل من دقيق غير الأرز قبل وضعه في سلة التخمير. والآن إذا بدأ العجين بالالتصاق بسلة التخمير بينما أخرجه ألجأ إلى إجراء جذرّي. أدهن فورًا قالب خبز بالزيت وأضع العجين داخله مباشرة. ويوفر قالب الخبز حاجزًا ولا يستطيع العجين الانبطح بقدرٍ ما. ولن يكون العجين منفّافًا بذلك الحد لكنه سيكون لذيذًا جدًّا إذا كنت تحب الخبز الحامض اللاذع.

وإذا كنت تملك مقياس حموضة \textenglish{(pH)} فدوّن ملاحظاتك عن درجة حموضة عجينك قبل الخَبز. سيسمح لك هذا بالحكم على عجينك أفضلًا طوال عملية التخمير.

\subsection{خبزي ينبطح أثناء التشكيل}

شأنه شان العجين الذي ينبطح بعد إخراجه من سلة التخمير، فالعجين المنبطح بعد التشكيل هو أيضًا إشارة محتملة إلى التخمير الزائد.

وعندما تحاول تشكيل العجين، هل يمكنك بسهولة أن تمزّق قطعًا من العجين؟ إذا كان كذلك فقد خمّرت عجينك زيادةً بلا شك. وإذا لم يكن كذلك فقد يكون مجرد إشارة إلى أنك لم تصنع ما يكفي من قوة العجين. فالتشباتا مثلًا هي عجين يميل إلى الانبطح قليلًا بعد التشكيل.

وإذا لم يكن عجينك قادرًا على التشكيل أصلًا فاستخدم قالب خبز مدهون لإنقاذ عجينك. يمكنك أيضًا أن تقطع قطعة من العجين وتستخدمها بادئةً لعجينك القادم. فعجينك الحامض في جوهره مجرد بادئة عملاقة.

\subsection{قشرتي تصبح مطاطية}

بحسب نمط الخبز الذي تصنعه فالقشرة السميكة المتشققة مطلوبة أحيانًا. وتُخلَق قشرة خبزك خلال المرحلة الثانية من عملية الخَبز بعد إزالة مصدر البخار من فرنك. والألوان الداكنة تُخلَق بالعملية المعروفة باسم تفاعل مايار ثم تتبعها عملية أخرى تُعرف بالتكرمل. وكل لون من القشرة يقدّم للتذوّق عِطرًا مختلفًا.

وما يحدث كثيرًا هو أن القشرة تصبح مطاطية بعد يوم. وأحيانًا عند الخَبز في المناطق المدارية ذات الرطوبة العالية لا تبقى القشرة في هذه المرحلة سوى بضع ساعات. وبعدها تصبح القشرة مطاطية. فهي لم تعد مقرمشة كالمقارنة مع اللحظة بعد الخَبز. ما يزال عجينك يحتوي رطوبةً. وستبدأ هذه الرطوبة في التجانس داخل الخبز النهائي وتبخر جزئيًا. والنتيجة أن قشرتك تصبح مطاطية.

وبالمثل، عند تخزين خبزك في وعاء أو في كيس بلاستيكي، ستصبح قشرتك مطاطية. لا يوجد لدي حل لهذا بعد. أميل عادةً إلى تخزين أخبازي في كيس بلاستيكي داخل ثلاجتي. وهذا يسمح للرطوبة بالبقاء داخل الخبز. وعند أخذ شريحة أحمّص كل شريحة دائمًا. وبهذا يعود بعض القرمشة. وإذا كنت تعرف طريقة عظيمة لذلك فتواصل معي وسأحدّث هذا الكتاب بنتائجك.

% Unit 15 continues here: \input{troubleshooting/crumb-structures}
% territory (accepted p. 166 ff.) — the 14.4 section head and intro
% paragraph above close accepted p. 165 and are included so the unit
% build carries the full accepted p. 165 content. Build-5 boundary
% note: the ~6 pt island-line drift (Latin-island metrics residual,
% R3 agenda) lets the 14.3.5 heading sit at the bottom of the unit
% build's p. 18 where the accepted edition opens p. 165 with it, so
% the unit build's p. 19 carries 14.3.5's body plus the 14.4 head and
% intro paragraph (documented deviation, unit 14 report).

\section{تشخيص بنية اللُّبّ}%
\label{sec:debugging-crumb-structure}

تمنحك بنية اللُّبّ في خبزك مؤشرات عن مدى نجاح عملية التخمير التي مرّت بها. ويمكنك أيضًا أن تكتشف العيوب الشائعة الناجمة عن تقنية غير سليمة. وسيمدّك هذا الفصل بمعلومات يمكنك استخدامها لتشخيص عملية الخَبز عندك.

% ---------------------------------------------------------------------------
% Unit 15: 14.4 body (crumb-structures) + 14.5 (misc tail), accepted
% pp. 166-179. Text layer transcribed from the accepted edition (PRIMARY);
% recovery cross-check and spacing/wording arbiter; upstream
% crumb-structures.tex + misc.tex tail give structure only (labels, cite
% keys, figure specs, plot includes). Figure placement [!htb] throughout,
% per the accepted edition's page-top figure placement (unit 14 part 1
% precedent); upstream [h]/[ht] specs not carried.
% ---------------------------------------------------------------------------

\begin{figure}[!htb]
  \centering
  \includegraphics[width=0.75\textwidth]{crumb-structures-book}
  \caption[تشخيص بنية اللُّبّ] {رسم تخطيطي يصوّر هياكل لُبٍّ مختلفة وأسبابها المقابلة. فلب الرغيف النهائي مفتاح أساسي لتحديد المشكلات المحتملة المتعلقة بالتخمير أو بتقنية الخَبز.}%
  \label{fig:crumb-structures-book}
\end{figure}

\subsection{التخمير المثالي}

\begin{figure}[!htb]
  \centering
  \includegraphics[width=\textwidth]{open-crumb}
  \caption[خبز مخمَّر جيدًا] {للخبز لُبٌّ منفتح بعض الشيء مع مناطق ذات بنية خلية النحل.}%
  \label{fig:open-crumb}
\end{figure}

قد يكون التخمير المثالي موضوعًا مثيرًا للنقاش وشديد الذاتية. فبالنسبة إليّ فإن خبز البادئ الحامضي المثالي يملك قشرة مقرمشة مقرونة بلُبٍّ رقيق منفتح بعض الشيء. وهذا هو التوازن المثالي بين القوامات المختلفة حين تقضم قضمة.

وهناك من يطاردون لُبًّا منفتحًا جدًّا، أي ما يعني وجود حويصلات كبيرة. ومن الذاتية أن تحدّد إن كان هذا أسلوب الخبز الذي يعجبك؛ غير أن بلوغه يتطلب أن تخمّر عجين خبزك تخميرًا مثاليًا. وبلوغ بنية لُبٍّ كهذه يلزمه قدر كبير من المهارة، من حيث إتقان التخمير ومن حيث التقنية معًا.

أنا شخصيًا أحب الخبز المنفتح، غير أن لُبّي المفضل أقل وحشيةً بعض الشيء. وأسلوب اللُّبّ الذي يعجبني يُسمى بنية خلية النحل \textenglish{(honeycomb crumb)}. فهو ليس منفتحًا أكثر من اللازم، بل منفتح في القدر الذي يجعل الخبز رقيقًا جدًّا. وبلوغ اللُّبّ المنفتح المذكور آنفًا يقتضي أن تلمس عجينك بأقل قدر ممكن. فكلما تفاعلت مع عجينك أكثر فرّغته من غازه أكثر. والإفراط في لمس العجين يؤدي إلى اندماج حويصلاته مع بعضها، فلن يكون اللُّبّ منفتحًا بالقدر نفسه. ولهذا فإن بلوغ لُبٍّ كهذا ينجح على أفضل وجه إن خمّرت رغيفًا واحدًا فقط في كل مرة. فعادةً إن اضطررت إلى تشكيل عجينك تشكيلًا أوليًا فستفرّغ عجينك من غازه قليلًا تلقائيًا أثناء التدوير. وإن تخطّيت هذه الخطوة وشكّلت عجينك مباشرةً فستبلغ لُبًّا أكثر انفتاحًا. والقاعدة الذهبية الجيدة هي ألا تلمس عجينك لمدة \LR{\textenglish{1\textendash 2}} ساعة على الأقل قبل التشكيل، لتبلغ أقصى انفتاح ممكن للُّبّ.

\begin{figure}[!htb]
  \centering
  \includegraphics[width=\textwidth]{honeycomb}
  \caption[بنية خلية النحل] {خبز حامض من دقيق كامل تقريبًا كلُّبّه من بنية خلية النحل.}%
  \label{fig:honeycomb}
\end{figure}

الآن يصبح هذا إشكاليًّا حين تريد صنع عدّة أرغفة في الوقت نفسه. فالتشكيل الأولي أمر جوهري إذ يلزمك تقسيم عجينك الأول الكبير إلى أجزاء أصغر. ومن دون عملية التشكيل الأولي سينتهي بك الأمر إلى عجائن خبز غير متجانسة عدّة. وتُستخدم هذه التقنية أيضًا عند صنع التشباتا؛ فهي عادةً لا تُشكَّل، وإنما تُقطع كتلة العجين فقط إلى قطع أصغر محاولًا العمل على العجين بأقل قدر ممكن. وبالتشكيل الأولي ستقترب حويصلات عجينك من بنية خلية النحل أكثر، إذ تندمج الجيوب الهوائية الكبيرة بعضها مع بعض قليلًا. وكما في بنية اللُّبّ المنفتح، يلزمك أيضًا هنا أن تُتقن عملية التخمير تمامًا لتبلغ هذا اللُّبّ. فالتخمير الأطول من اللازم سيجعل الغاز يتسرب من عجينك أثناء الخَبز، فلا تستطيع أقراص خلية النحل الاحتفاظ بالغاز. وإن خمّرت مدةً أقصر من اللازم فلن يكون هناك غاز كافٍ لنفخ الهياكل. وبالنسبة إليّ هذا هو أسلوب اللُّبّ المثالي. وكما أني من محبي المربى، فلن يسقط أي مربى من شريحة من هذا الخبز بخلاف ما يحدث مع اللُّبّ المنفتح.

\subsection{المُخمَّر بإفراط}%
\label{subsec:overfermented-dough}

\begin{figure}[!htb]
  \centering
  \includegraphics[width=\textwidth]{fermented-too-long}
  \caption[خبز حامض مخمَّر بإفراط] {عجين منبطح نسبيًا يضم جيوبًا هوائية صغيرة عدّة.}%
  \label{fig:fermented-too-long}
\end{figure}

عند تخمير عجينك لمدةً أطول من اللازم، يبدأ أنزيم البروتييز بتفكيك غلوتين دقيقك. علاوة على ذلك تستهلك البكتيريا الغلوتين في عملية تُسمى تفكك البروتين~\cite{raffaella+di+cagno}. ويشير الخبّازون إلى هذه العملية أيضًا باسم تعفّن الغلوتين \textenglish{(gluten rot)}. والغلوتين الذي يحبس عادةً غاز \LR{\ch{CO2}} الناتج عن تخمير الكائنات الدقيقة لديك لم يعد قادرًا على إبقاء الغاز داخل العجين. فيتبدد الغاز نحو الخارج فينتج عن ذلك حويصلات أصغر في لُبّك. ويميل الخبز نفسه إلى الانبطح الشديد في الفرن. ويشير الخبّازون غالبًا إلى هذا الأسلوب من الخبز بوصفه قرصًا مسطّحًا \textenglish{\emph{pancake}}. ويمكن مقارنة انتفاخ الفرن فيه بعجائن الخبز المصنوعة من دقيق قليل الغلوتين كالقمح الأحادي.

سيتميز خبزك بحموضة كبيرة وبطعم حامض قوي مميز. ومن منظور الذوق، قد يكون حامضيًا أكثر من اللازم بعض الشيء. ومن تجاربي الخاصة مع العائلة والأصدقاء \textenglish{(n=15\textendash 20)} أستطيع أن أقول إن هذا الأسلوب من الخبز يُقدَّر عادةً بدرجة أقل. غير أني أحب شخصيًا الطعم القوي الدسم هذا. فهو ممتاز بقرين شيء حلو أو مع حساء. ومن منظور القوام، لم يعد رقيقًا بالقدر الذي يمكن أن يكون عليه. وقد يتكون في اللُّبّ ما يشبه الجلد. والسبب أن بنية الغلوتين تفككت بفعل البكتيريا تفككًا كبيرًا. علاوة على ذلك يملك هذا الأسلوب من الخبز كميةً من الغلوتين أقل بكثير~\cite{raffaella+di+cagno} ولم يعد يُقارن بالدقيق الخام؛ إنه منتج مخمَّر تمامًا. ويمكنك مقارنته بجبن أزرق خالٍ من اللاكتوز تقريبًا.

وحين تحاول العمل على العجين ستلاحظ فجأةً أنه صار لزجًا جدًّا. فلم تعد تستطيع تشكيل العجين والعمل عليه على النحو الصحيح. وعند محاولة إخراج العجين من سلة التخمير سينبطح العجين انبطاحًا كبيرًا. وعلاوة على ذلك قد يلتصق عجينك بسلة التخمير في حالات عدّة. وعند بداياتي في الخَبز كنت أستخدم الكثير من دقيق الأرز في سلة التخمير كي أجفف سطح العجين كثيرًا. وبهذه الطريقة كان العجين لا يلتصق رغم أنه مخمَّر بإفراط. غير أن الأمر انكشف أن مشكلة اللزوجة كانت ناتجة عن قصور فهمي لعملية التخمير. والآن لا أستخدم دقيق الأرز إطلاقًا، إلا حين أحاول عمل تشريط زخرفي.

وإذا لاحظت أثناء تطويات وفرُد أو أثناء التشكيل أن عجينك صار فجأةً لزجًا بشكل مفرط، فأفضل خيار هو استخدام قالب الخبز. خذ عجينك ببساطة وألقِ به في قالب الخبز. وانتظر حتى تزداد كتلة العجين في الحجم بعض الشيء من جديد ثم اخبزها. وسيكون عندك خبز حامض طيب الطعم جدًّا. وإن كان حامضيًا أكثر من اللازم بعض الشيء، فيمكنك ببساطة أن تخبز عجينك مدةً أطول لكي يتبخر بعض الحموضة أثناء عملية الخَبز. ويمكنك أيضًا أن تستخدم عجينك لإعداد بادئة حامضية جديدة وتحاول من جديد غدًا. وأخيرًا، إن كنت جائعًا فيمكنك ببساطة أن تصب بعض عجينك مباشرةً في مقلاة مسخَّنة مع قليل من الزيت. فسيصنع لك أقراصًا مسطّحة لذيذة من الخبز الحامض.

ولإصلاح المشكلات المتعلقة بالتخمير الزائد، عليك أن توقف عملية التخمير مبكرًا. وما أحب فعله هو استخراج عيّنة تخمير صغيرة من عجيني. وبحسب الزيادة في حجم هذه العيّنة أستطيع غالبًا أن أحكم متى يكتمل تخميري. حاول أن تبدأ بزيادة حجم تبلغ \LR{\qty{25}{\percent}} من عجينك الرئيسي أو عيّنتك. وبحسب كمية الغلوتين في دقيقك يمكنك التخمير مدةً أطول. فباستخدام دقيق قوي تبلغ نسبة البروتين فيه \LR{\qtyrange{14}{15}{\percent}} ينبغي أن تستطيع التخمير بأمان حتى زيادة حجم تبلغ \LR{\qty{100}{\percent}}. لكن هذا يعتمد أيضًا على تركيبة بادئتك الحامضية من خميرة وبكتيريا. فكلما زاد التخمير البكتيري تفككت بنية عجينك أسرع. والتغذية المتكررة لبادئتك الحامضية ستسهم في تحسين نشاط الخميرة. علاوة على ذلك قد تكون البادئة الحامضية الصلبة حلًّا جيدًا أيضًا. فنشاط الخميرة المعزز سيؤدي إلى عجين أكثر رقة ونشاطًا بكتيريًا أقل. ونشاط الخميرة الأفضل سيسهم أيضًا في حموضة أقل في خبزك النهائي. وإن كنت من طاردَي النكهة الحامضة القوية جدًّا فسيساعدك دقيق أقوى يملك غلوتينًا أكثر.

وعند إبطاء تخمير الخبز الحامض (التخمير البارد في الثلاجة)، تلعب درجة الحرارة دورًا محوريًا في سرعات التخمير. فكلما برد العجين في الثلاجة تباطأ التخمير. وبدء عملية الإبطاء من حرارة أدفأ يعني أن هذا التباطؤ يستغرق وقتًا أطول. فمثلًا، عجين يصبح مثاليًا بعد \LR{8} ساعات من الإبطاء قد يجهز في \LR{4} ساعات فقط إن بدأ من حرارة أعلى. ولهذا فإن إجراء التجارب وتحديد مدة الإبطاء المثالية لظروفك الخاصة أمر بالغ الأهمية. وبالعكس، إن بدأ العجين أبرد فسيتوقف التخمير أسرع في الثلاجة. وفي هذه السيناريوهات قد يكون من المفيد السماح للعجين بالإثمار في حرارة الغرفة باختصار قبل وضعه في الثلاجة.

\subsection{المُخمَّر بأقل من اللازم}

\begin{figure}[!htb]
  \centering
  \includegraphics[width=\textwidth]{fermented-too-short-underbaked}
  \caption[خبز مخمَّر بأقل من اللازم] {عجين كثيف يضم منطقة لزجة لم تدخل الجيلاتنة تمامًا. وقد أُرسل الصورة المستخدم \textenglish{\emph{wahlfeld}} من خادم دِسكورد \textenglish{(Discord)} التابع لمجتمعنا.}%
  \label{fig:fermented-too-short-underbaked}
\end{figure}

يُشار إلى هذا العيب شائعًا أيضًا باسم قلة الإثمار \textenglish{(underproofed)}. غير أن قلة الإثمار ليس مصطلحًا جيدًا لأنه يشير فقط إلى قِصَر مرحلة التخمير النهائي في عملية صناعة الخبز. ولو خبزت خبزك بعد مرحلة تخمير أول مثالية التوقيت، فلن تُوصف النتيجة بقلة الإثمار حتى لو تخطّيت مرحلة التخمير النهائي بأكملها. فالإثمار سيجعل عجينك قابلًا للمد بعض الشيء ويتيح لبادئك الحامضي أن ينفخ العجين أكثر قليلًا. وحين تواجه خبزًا مخمَّرًا بأقل من اللازم، فإن شيئًا ما قد أخفق قبل ذلك أثناء مرحلة التخمير الأول، أو ربما قبل ذلك حتى في بادئتك الحامضية.

العجين النموذجي المخمَّر بأقل من اللازم يملك حويصلات كبيرة جدًّا وهو مبتل ولزج في بعض مناطق العجين. ويمكن مقارنة الحويصلات الكبيرة بصنع تورتيلا من القمح أو الذرة غير مخمَّرة. فحين تخبز العجين في المقلاة، يبدأ الماء بالتبخر ببطء. ويحتبس الغاز في بنية العجين فيخلق جيوبًا. وفي حالة التورتيلا فهذا هو السلوك المنشود. لكن حين تلاحظ هذه العملية في عجين أكبر، فستخلق عدّة حويصلات فائقة. يتبخر الماء وتتكون الحويصلات الأولى. ثم في لحظة ما يبلغ النشا الجيلاتنة ويتصلب. ويحدث هذا أولًا داخل الحويصلات إذ تسخن الأجزاء الداخلية أسرع مقارنةً ببقية العجين. وبمجرد أن يدخل النشأ كله طور الجيلاتنة تحتفظ الحويصلات بشكلها ولم تعد تتوسع. وأثناء هذه العملية تُدفع أجزاء أخرى من عجين الخبز نحو الخارج. ولهذا قد يتميز العجين المخمَّر بأقل من اللازم أحيانًا بالأذن حتى أثناء عملية الخَبز. ويُشار إلى هذا أيضًا شائعًا باسم لُبّ الأحمق. فأنت متحمس للأذن التي قد يصعب بلوغها جدًّا. وقد تظن أنك صنفت أخيرًا بعض الحويصلات الكبيرة في لُبّك. لكن في الواقع إنما خمّرت لمدةً أقصر من اللازم.

\begin{figure}[!htb]
  \centering
  \includegraphics[width=\textwidth]{fools-crumb}
  \caption[حويصلات كبيرة في لُبّ الأحمق] {مثال نموذجي على لُبّ الأحمق يتميز بأذن وعدّة حويصلات كبيرة أكثر من اللازم. وقد أُرسلت الصورة روشيل من خادم دِسكورد التابع لمجتمعنا.}%
  \label{fig:fools-crumb}
\end{figure}

في العجين المخمَّر على نحو سليم تسهم الحويصلات في انتقال الحرارة عبر العجين كله. فمن داخل الحويصلات الدقيقة العديدة الناتجة عن التخمير يدخل النشأ طور الجيلاتنة. أما في العجين المخمَّر بأقل من اللازم فهذا الانتقال الحراري لا يعمل على النحو الصحيح. ولهذا قد تجد أحيانًا مناطق تبدو كعجين خام. ويشير الخبّازون إلى ذلك أحيانًا بوصفه بنية لزجة جدًّا. وإن خبزت عجينك مدةً أطول لدخل النشأ في هذه المناطق طور الجيلاتنة أيضًا على نحو سليم. غير أن أجزاء أخرى من خبزك قد تُخبز حينها مدةً أطول من اللازم.

ولإصلاح المشكلات المتعلقة بالنقص في التخمير، عليك ببساطة أن تخمّر عجينك مدةً أطول. والآن، هناك حد أعلى لزمن التخمير إذ يبدأ دقيقك بالتفكك لحظة ملامسته الماء. ولهذا قد تكون فكرة جيدة أن تسرّع ببساطة عملية التخمير لديك. وبصورة تقريبية أحاول أن أستهدف زمن تخمير أول يبلغ نحو \LR{\textenglish{8\textendash 12}} ساعة عادةً. وبلوغ ذلك يمكنك أن تحاول أن تجعل بادئتك الحامضية أكثر نشاطًا. ويمكن فعل ذلك بتغذية بادئتك يوميًا على مدى عدّة أيام. استخدم النسبة نفسها التي تستخدمها لعجين خبزك الرئيسي. وبافتراض أنك تستخدم \LR{\qty{20}{\percent}} بادئة محسوبةً على الدقيق، فاستخدم نسبة \LR{1:5:5} لتغذية بادئتك. وهذه تعني مثلًا \LR{\qty{10}{\gram}} من البادئة الموجودة، و\LR{\qty{50}{\gram}} من الدقيق، و\LR{\qty{50}{\gram}} من الماء. ولزيادة نشاط خميرتك أكثر فيمكنك أن تفكر في صنع بادئة حامضية صلبة. فالبكتيريا تنتج الحمض في الغالب. وكلما تراكمت الحموضة قلّ نشاط خميرتك. وتتيح لك البادئة الحامضية الصلبة أن تبدأ تخمير عجينك بنشاط خميرة أقوى ونشاط بكتيري أقل.

\subsection{قوة العجين غير الكافية}

\begin{figure}[!htb]
  \centering
  \includegraphics[width=\textwidth]{flat-bread}
  \caption[خبز منبطح جدًّا بسبب قوة عجين غير كافية] {خبز منبطح جدًّا بسبب قوة عجين غير كافية.}%
  \label{fig:flat-bread}
\end{figure}

وحين ينبطح العجين انبطاحًا كبيرًا جدًّا أثناء عملية الخَبز، فالأرجح أنك لم تصنع قوة عجين كافية. ويعني ذلك أن مصفوفة الغلوتين لديك لم تتطور على النحو الصحيح. فالعجين قابل للمد أكثر من اللازم وينبطح في الغالب بدل أن ينتفخ نحو الأعلى في الفرن. وقد يحدث هذا أيضًا إن أثمرت عجينك مدةً أطول من اللازم. فمع الزمن يرتخي الغلوتين ويصبح عجينك قابلًا للمد أكثر فأكثر. ويمكنك أيضًا أن تلاحظ سلوك ارتخاء الغلوتين عند صنع فطيرة بيتزا. فبعد تشكيل كرات العجين مباشرةً من الصعب جدًّا شدّ فطيرة البيتزا إلى شكلها. وإن انتظرت \LR{\textenglish{30\textendash 90}} دقيقة صار شدّ العجين أسهل بكثير.

أيسر طرق إصلاح ذلك على الأرجح أن تعجن عجينك أكثر في البداية. ولبساطة الأمور فكّر في استخدام ماء أقل لدقيقك أيضًا. وستنتج عن ذلك عجينة أكثر مرونة على الفور. وهذا مفهوم شائع الاستخدام في أسلوب الخبز الحامض بلا عجن \textenglish{no-knead}. ويمكنك بدلًا من ذلك أيضًا أن تجري تطويات وفرُد أكثر أثناء عملية التخمير الأول. فكل تطوية وفرُد ستسهم في تقوية مصفوفة الغلوتين وصنع عجينة أكثر مرونة. والخيار الأخير لإصلاح عجين ذي قوة غير كافية أن تشكّل عجينك تشكيلًا أكثر إحكامًا.

\subsection{خُبز على حرارة مرتفعة جدًّا}

\begin{figure}[!htb]
  \centering
  \includegraphics[width=\textwidth]{baked-too-hot-v2}
  \caption[خبز بحويصلات كبيرة جدًّا قرب القشرة] {خبز بحويصلات كبيرة جدًّا قرب القشرة.}%
  \label{fig:baked-too-hot}
\end{figure}

هذا خطأ شائع حدث لي كثيرًا. فحين تخبز عجينك على درجة حرارة مرتفعة أكثر من اللازم فإنك تقيّد تمدد عجينك. فيدخل النشا طور الجيلاتنة ويصبح صلبًا أكثر فأكثر. وبمجرد أن يجف السطح وبلغ نحو \qtyrange{105}{120}{\degreeCelsius} (\qtyrange{221}{248}{\degF}) يبدأ تفاعل مايار بتسميك قشرة عجين خبزك تسميكًا تامًّا~\cite{purlis+browning}. وهذا شبيه بخبز عجين خبزك من دون بخار. وبما أن الحرارة الداخلية للعجين ترتفع، يتبخر ماء أكثر فأكثر ويتوسع الغاز ويُدفع العجين نحو الأعلى. وحالما يبلغ العجين القشرة لم يعد يستطيع التمدد. فتندمج الحويصلات في هياكل أكبر قرب سطح العجين. وبالخبز على حرارة مرتفعة جدًّا لن تبني الأذن التي تضيف نكهة إضافية. علاوة على ذلك، بتقييد تمدده لن يكون اللُّبّ رقيقًا بالقدر الذي يمكن أن يكون عليه.

وإن كان لديك عجين قابل للمد وعالي الترطيب، فإن الخبز على حرارة منخفضة جدًّا سيجعل العجين ينبطح انبطاحًا كبيرًا. فدخول النشا طور الجيلاتنة ضروري لكي يحتفظ العجين ببنية. وبعد إجراء عدّة تجارب، يبدو أن نقطة الاعتدال المثالية عندي لأقصى انتفاخ في الفرن تقع عند نحو \LR{\qty{230}{\degreeCelsius}} (\LR{\qty{446}{\degF}}). اختبر حرارة فرنك، لأنه في حالات عدّة قد لا تطابق الحرارة المعروضة الحرارة الفعلية لفرنك~\cite{too+hot+baking}. وتأكد من إطفاء مروحة فرنك. فمعظم أفران المنازل مصممة لطرد البخار بأسرع ما يمكن. وإن لم تستطع إطفاء المروحة ففكر في استخدام وعاء خبز هولندي.

\subsection{خُبز ببخار أقل من اللازم}

\begin{figure}[!htb]
  \centering
  \includegraphics[width=\textwidth]{no-steam}
  \caption[خبز مخبوز ببخار أقل من اللازم] {أحد أخبازي الأولى التي خبزتها في بيت صديق حيث لم أستطع تبخير العجين على نحو سليم.}%
  \label{fig:no-steam}
\end{figure}

مثل الخبز على حرارة مرتفعة جدًّا، حين تخبز من دون بخار كافٍ تتكون قشرة عجينك بسرعة كبيرة. ومن الصعب تمييز الفرق بين الخطأين. وعادةً أسأل أولًا عن درجة الحرارة ثم عن تقنية التبخير لأحدد ما الذي قد يكون خاطئًا في عملية الخَبز. ويمكن عادةً اكتشاف قلة البخار بوجود قشرة سميكة من كل الجهات حول عجينك مقرونة بحويصلات كبيرة باتجاه الحواف.

يمنع البخار في جوهره تفاعل مايار من الحدوث بسرعة كبيرة على قشرتك. ولهذا فالتبخير خلال المراحل الأولى من الخَبز مهم للغاية. فالبخار يُبقي حرارة قشرتك قريبة من نحو \LR{\qty{100}{\degreeCelsius}} (\LR{\qty{212}{\degF}}). وبلوغ البخار ممكن باستخدام وعاء خبز هولندي، أو صينية مقلوبة، أو وعاء من الماء المغلي، أو باختلافات بينها. وقد يكون عندك فرن مزوّد بوظيفة بخار مدمجة. وكل الطرق تنجح، والأمر يعتمد على ما هو بين يديك. وطريقتي الافتراضية المفضلة هي صينية مقلوبة فوق عجيني مقرونة بوعاء مليء بالماء المغلي باتجاه أسفل الفرن.

\begin{figure}[!htb]
  \centering
  \includegraphics[width=0.7\textwidth]{apple-experiment-temperatures}
  \caption[قياس الحرارة المحيطة وحرارة السطح] {تفاحة مزوّدة بمجسَّين لقياس الحرارة المحيطة وحرارة السطح لعدّة تقنيات تبخير في وعاء خبز هولندي.}%
  \label{fig:apple-experiment-temperatures}
\end{figure}

الآن يمكن أيضًا أن يكون هناك بخار أكثر من اللازم. ولأجل هذا اختبرتُ استخدام وعاء خبز هولندي مقرونًا بمكعبات ثلج كبيرة لتوفير بخار إضافي. فحرارة سطح عجيني ستقفز مباشرةً إلى ما يقارب \textenglish{100°C}. فالبخار يحوي طاقة أكبر وبالتالي يمكنه عبر الحمل الحراري أن يسخّن سطح عجينك أسرع. واختبرتُ ذلك بوضع تفاحة داخل وعاء خبز هولندي وقياس حرارة سطحها باستخدام مقياس حرارة الشوّاء. ثم غيّرتُ طرق التبخير لأرسم كيف تتغير الحرارة القريبة من السطح بسرعة. اختبرتُ مكعب ثلج داخل وعاء خبز هولندي مسخَّن مسبقًا، ووعاء خبز هولندي مسخَّنًا مسبقًا مجردًا، ووعاء خبز هولندي مسخَّنًا مسبقًا مع رشات ماء على سطح التفاحة، ووعاء خبز هولندي غير مسخَّن مسبقًا سخّنتُ جزءَه السفلي فقط. وقد أظهرت التجربة حينها أن طريقة مكعب الثلج تسخّن سطح التفاحة أسرع بكثير. وحين أعدتُ هذا بعجين خبز حققت انتفاخًا أقل في الفرن.

\begin{figure}[!htb]
  \centering
  \input{plots/fig-temperature-surface.tex}
  \caption[حرارة السطح مقابل تقنية التبخير] {مخطط يوضح كيف تتغير حرارة سطح التفاحة مع تقنيات التبخير المختلفة.}%
  \label{fig:apple-experiment-surface-temperatures}
\end{figure}

\begin{figure}[!htb]
  \centering
  \input{plots/fig-temperature-ambient.tex}
  \caption[حرارة الوعاء الهولندي مقابل تقنية التبخير] {يوضح هذا الشكل كيف تتغير الحرارات المحيطة داخل الوعاء الهولندي بحسب تقنية التبخير المستخدمة.}%
  \label{fig:apple-experiment-ambient-temperatures}
\end{figure}

عمومًا بلوغ بخار أكثر من اللازم أمر صعب نسبيًا. فلم أستطع ارتكاب هذا الخطأ إلا عند استخدام وعاء الخبز الهولندي بوصفه طريقة التبخير مقرونًا بمكعبات ثلج كبيرة نسبيًا. وبعد الحديث مع خبازين آخرين يستخدمون الوعاء الهولندي نفسه، يبدو أن مكعبات الثلج عندي (نحو \LR{\qty{80}{\gram}}) كانت أثقل بأربع مرات من التي يستخدمها الخبازون الآخرون (\LR{\qty{20}{\gram}}).

\section{متفرقات}

\subsection{الخَبز في المناطق المدارية}

بحسب الحرارة يتأقلم سرعة تخميرك. ففي بيئة أدفأ كل شيء أسرع. وفي بيئة أبرد كل شيء أبطأ.

وهذا يشمل السرعة التي تخمّر بها البادئة الحامضية العجين، بل أيضًا سرعة التفاعلات الأنزيمية. فأنزيمات الأميلاز والبروتييز تعمل أسرع، فتجعل المزيد من السكريات متاحًا وتفكك بروتينات الغلوتين.

عند نحو \LR{\qty{22}{\degreeCelsius}} (\LR{\qty{72}{\degF}}) في مطبخي يكون تخميري الأول جاهزًا بعد نحو \LR{10} ساعات. أستخدم نحو \LR{\qty{20}{\percent}} من بادئة حامضية محسوبةً على الدقيق. وفي الوقت الصيفي ترتفع حرارة مطبخي أحيانًا إلى \LR{\qty{25}{\degreeCelsius}} (\LR{\qty{77}{\degF}}). وفي تلك الحالة أخفّض البادئة الحامضية إلى نحو \LR{\qty{10}{\percent}}.

ولو لم أفعل ذلك لكان تخميري قد اكتمل بعد نحو \LR{\textenglish{4\textendash 7}} ساعات. والمشكلة أن العجين غير مستقر بما يكفي عند التخمير بهذه السرعة العالية. ومعنى ذلك أنك تصطدم بسهولة بمشكلات التخمير الزائد. ويصبح بلوغ نقطة الاعتدال المثالية بين التخمير الكافي والتخمير الزائد أصعب بكثير. وعادةً قد يكون عندك نافذة زمنية تبلغ ساعة واحدة. لكنها عند السرعة السريعة قد تنكمش إلى نافذة زمنية تبلغ \LR{20} دقيقة. والآن عند \LR{\qty{30}{\degreeCelsius}} (\LR{\qty{86}{\degF}}) يتحرك كل شيء أسرع بكثير. فقد يكون تخميرك الأول مكتملًا في \LR{\textenglish{2\textendash 4}} ساعات عند استخدام بادئة \LR{\qtyrange{10}{20}{\percent}}. ويصبح تخمير عجينك النهائي في الثلاجة شبه مستحيل. فكلما برد عجينك في الثلاجة تباطأ التخمير أيضًا. لكن تبريد العجين من \LR{\qty{30}{\degreeCelsius}} إلى \LR{\qtyrange{4}{6}{\degreeCelsius}} في ثلاجتك يستغرق مدةً أطول بكثير. فعجينك أكثر نشاطًا بكثير مقارنةً بعجين يبدأ بحرارة تبلغ \LR{\qtyrange{20}{25}{\degreeCelsius}}. وقد ينتهي بك الأمر إلى تخمير نهائي زائد إذا تركته طوال الليل في الثلاجة.

ولهذا السبب أنصح بأن تخفّض كمية البادئة التي تستخدمها في المناطق المدارية إلى نحو \LR{\qtyrange{1}{5}{\percent}} محسوبةً على الدقيق. سيعمل هذا على إبطاء عملية التخمير إبطاءً كبيرًا ويمنحك نافذة زمنية أكبر. حاول أن تستهدف تخميرًا أولًا كليًّا لا يقل عن \LR{\textenglish{8\textendash 10}} ساعات. خفّض كمية البادئة للوصول إلى هناك.

وعند صنع العجين حاول استخدام درجة حرارة الماء نفسها كحرارتك المحيطة. وبافتراض أن الحرارة ستتسلق إلى \LR{\qty{30}{\degreeCelsius}} فحاول بدء عجينك بماء \LR{\qty{30}{\degreeCelsius}}. ويعني ذلك أنك تستطيع الاعتماد بحذرٍ صغير على عيّنة تخمير صغيرة (جرّة أليكوت) تُجمّل لك سير تخميرك. ولمزيد من القراءة عن هذه التقنية راجع القسم \LR{9.9}.

ولا تعمل العيّنة بموثوقية إلا إذا كانت حرارة عجينك مساويةً لحرارتك المحيطة. وإلا فإن العيّنة تسخن أو تبرد أسرع. فتحرّ بحذر عند استخدام العيّنة في هذه الحالة. ومن الأفضل دائمًا أن توقف التخمير قبل قليلٍ من أن تكون متأخرًا. وسيساعدك أثناء التخمير الأول على تطوير إحساسٍ أفضل بالعجين. والأداة المكلفة لكنها قد تكون مفيدة هي مقياس حموضة \textenglish{(pH)} يتيح لك قياس مقدار الحموضة التي خلقتها بكتيريا حمض اللبنيك وحمض الخليك قياسًا مثاليًا. وفي هذه الحالة قس درجة الحموضة مرارًا وتكرارًا وافهم قيمةً تناسب بادئتك الحامضية. وفي حالتي أميل إلى إنهاء التخمير الأول عند درجة حموضة تبلغ نحو \LR{4.1}. رجاءً لا تتبع قيمة درجة الحموضة لي مجرد التباع؛ فهي فردية جدًّا. حافظ على القياس بعجائن مختلفة لتفيق قيمةً تناسبك.

\subsection{دقيق منخفض الغلوتين، فماذا أفعل؟}

يمكنك دائمًا أن تخلط قليلًا من الغلوتين الحيوي \textenglish{(vital wheat gluten)}. فالغلوتين الحيوي غلوتين مركّز مستخرج من دقيق القمح.

أنصحك بإضافة نحو \LR{\qty{5}{\gram}} من الغلوتين الحيوي لكل \LR{\qty{100}{\gram}} من الدقيق الذي تستخدمه.


\end{document}
